% =====================================================================
%  CARNET D'ENTRAÎNEMENT — FICHE 8 — THÈME 1
%  Tuto : Premier principe et enthalpie (exo/endothermique)
%  Compilation : pdflatex
% =====================================================================
\documentclass[11pt,a4paper]{article}
\usepackage[utf8]{inputenc}
\usepackage[T1]{fontenc}
\usepackage{lmodern}
\usepackage[french]{babel}
\usepackage{amsmath,amssymb,mathtools}
\usepackage{icomma}
\usepackage{siunitx}
\sisetup{output-decimal-marker={,},per-mode=power,group-separator={\,},group-minimum-digits=5}
\usepackage[version=4]{mhchem}
\usepackage{xcolor}
\usepackage{colortbl}
\usepackage{tikz}
\usepackage{pgfplots}
\usepackage[most]{tcolorbox}
\usepackage{enumitem}
\usepackage{array}
\usepackage{booktabs}
\usepackage{multicol}
\usepackage{fancyhdr}
\usepackage{caption}
\captionsetup{labelformat=empty,justification=centering,font=small}
\usepackage[hidelinks]{hyperref}
\usetikzlibrary{arrows.meta,positioning,calc,decorations.pathmorphing,decorations.markings,angles,quotes,patterns,backgrounds,shapes.geometric,shapes.misc,fadings,intersections,fit}
\pgfplotsset{compat=1.18}
\usepackage[a4paper,margin=2.1cm,headheight=26pt,headsep=10pt]{geometry}
\usepackage{titlesec}

% ---------- palette ----------
\definecolor{primary}{RGB}{150,98,162}
\definecolor{primarydark}{RGB}{92,52,110}
\definecolor{defcol}{RGB}{0,114,178}
\definecolor{thmcol}{RGB}{0,140,120}
\definecolor{formulacol}{RGB}{198,93,9}
\definecolor{remarkcol}{RGB}{120,94,200}
\definecolor{attncol}{RGB}{200,40,55}
\definecolor{excol}{RGB}{0,135,90}
\definecolor{methcol}{RGB}{90,90,100}
\definecolor{exocol}{RGB}{200,60,40}
\definecolor{endocol}{RGB}{0,110,180}
\definecolor{chaleur}{RGB}{214,72,33}
\definecolor{travail}{RGB}{30,110,180}

\tcbset{boxbase/.style={enhanced,breakable,boxrule=0.9pt,arc=2.5pt,left=3mm,right=3mm,top=2.3mm,bottom=2.3mm,fonttitle=\bfseries,coltitle=white,toptitle=0.6mm,bottomtitle=0.6mm}}
\newtcolorbox{definitionb}[1][]{boxbase,colback=defcol!5,colframe=defcol,title={\textsc{Définition}\ifx\relax#1\relax\relax\else\textmd{ --- \itshape #1}\fi}}
\newtcolorbox{proprieteb}[1][]{boxbase,colback=thmcol!6,colframe=thmcol,title={\textsc{Propriété}\ifx\relax#1\relax\relax\else\textmd{ --- \itshape #1}\fi}}
\newtcolorbox{formuleb}[1][]{boxbase,colback=formulacol!6,colframe=formulacol,title={\textsc{Formule}\ifx\relax#1\relax\relax\else\textmd{ --- \itshape #1}\fi}}
\newtcolorbox{methodeb}[1][]{boxbase,colback=methcol!7,colframe=methcol,sharp corners,title={\textsc{Méthode}\ifx\relax#1\relax\relax\else\textmd{ : \itshape #1}\fi}}
\newtcolorbox{remarqueb}[1][]{boxbase,colback=remarkcol!6,colframe=remarkcol,title={\textsc{Remarque}\ifx\relax#1\relax\relax\else\textmd{ : \itshape #1}\fi}}
\newtcolorbox{attentionb}[1][]{boxbase,colback=attncol!6,colframe=attncol,title={\textsc{Attention}\ifx\relax#1\relax\relax\else\textmd{ : \itshape #1}\fi}}
\newtcolorbox{exempleb}[1][]{boxbase,colback=excol!5,colframe=excol,title={\textsc{Exemple}\ifx\relax#1\relax\relax\else\textmd{ : \itshape #1}\fi}}

\newenvironment{ou}{\par\smallskip\noindent\textbf{où :}\nopagebreak\begin{itemize}[leftmargin=1.8em,topsep=2pt,itemsep=1.5pt,parsep=0pt]}{\end{itemize}}

\newcommand{\dH}{\Delta H}
\newcommand{\dHr}{\Delta H^{\circ}_{\text{r}}}
\newcommand{\dHf}{\Delta H^{\circ}_{\text{f}}}
\newcommand{\dS}{\Delta S}
\newcommand{\dU}{\Delta U}

\titleformat{\section}{\large\bfseries\color{primarydark}}{\thesection}{0.6em}{}
\titleformat{\subsection}{\normalsize\bfseries\color{primary}}{\thesubsection}{0.6em}{}

\pagestyle{fancy}
\fancyhf{}
\fancyhead[L]{\small\textcolor{primarydark}{\textbf{Fiche 8} — Thème 1 : Premier principe et enthalpie}}
\fancyhead[R]{\small\textcolor{primarydark}{\textsc{Tuto}}}
\fancyfoot[C]{\small\thepage}
\renewcommand{\headrulewidth}{0.8pt}
\renewcommand{\headrule}{\hbox to\headwidth{\color{primary}\leaders\hrule height \headrulewidth\hfill}}

\setlength{\parindent}{0pt}
\setlength{\parskip}{3pt}

\begin{document}

\begin{center}
\begin{tikzpicture}
\node[rectangle,rounded corners=7pt,fill=primary,text=white,minimum width=16cm,
      minimum height=1.5cm,align=center,inner sep=8pt]
  {{\Large\bfseries Thème 1 — Premier principe et enthalpie}\\[2pt]
   {\small Énergie interne, chaleur, travail, réactions exo\,/\,endothermiques}};
\end{tikzpicture}
\end{center}

\section{L'idée générale}
Une réaction chimique réorganise les liaisons : elle \textbf{échange de l'énergie} avec l'extérieur (chaleur
$Q$ et travail $W$). Ce thème répond à deux questions : \emph{de quel côté va l'énergie} (signe de $Q$, $W$, $\dH$) et \emph{combien}.

\begin{definitionb}[système ouvert / fermé / isolé]
Le \textbf{système} est la portion étudiée (le contenu du bécher) ; le reste est le \textbf{milieu
extérieur}. Un système est \textbf{ouvert} (échange matière \emph{et} énergie), \textbf{fermé} (énergie
seulement) ou \textbf{isolé} (ni l'un ni l'autre). Un calorimètre parfait est un système isolé, dit
\emph{adiabatique}.
\end{definitionb}

\section{Le premier principe et les conventions de signe}

\begin{formuleb}[Premier principe de la thermodynamique]
\[ \boxed{\;\dU \;=\; Q \;+\; W\;}\qquad\text{(énergies en }\si{\joule}\text{)} \]
\begin{ou}
  \item $\dU$ : variation d'\textbf{énergie interne} du système ;
  \item $Q$ : \textbf{chaleur} échangée ; \quad $W$ : \textbf{travail} échangé.
\end{ou}
\end{formuleb}

\begin{methodeb}[le réflexe du signe — point de vue du système]
Ce que le système \textbf{reçoit} est \textbf{positif} ; ce qu'il \textbf{cède} est \textbf{négatif}.
\begin{itemize}[leftmargin=1.5em,topsep=1pt,itemsep=1pt]
  \item chaleur reçue $\Rightarrow Q>0$ \; ; \; chaleur cédée $\Rightarrow Q<0$ ;
  \item travail reçu $\Rightarrow W>0$ \; ; \; travail cédé $\Rightarrow W<0$.
\end{itemize}
Système \textbf{isolé} : $Q=0$ et $W=0$, donc $\dU=0$ ($U$ constante).
\end{methodeb}

À \textbf{pression constante} (cas usuel : réaction à l'air libre), deux résultats servent tout le temps :
\begin{center}
\begin{tabular}{@{}p{7.3cm}p{8.3cm}@{}}
\textbf{Travail des forces de pression :} \quad $W = -\,p\,\Delta V$ &
\textbf{Chaleur à pression constante :} \quad $Q_p = \dH$ \\
($p$ en \si{\pascal}, $\Delta V=V_{\text{f}}-V_{\text{i}}$ en \si{\cubic\metre}) &
($\dH$ = variation d'\textbf{enthalpie}, en \si{\joule} ou \si{\kilo\joule}) \\
\end{tabular}
\end{center}
Le signe « $-$ » de $W=-p\Delta V$ traduit la convention : si le gaz se \emph{dilate} ($\Delta V>0$) il
\emph{pousse} l'extérieur, donc il \emph{cède} du travail ($W<0$).

\section{Réactions exothermiques et endothermiques}

Comme $Q_p=\dH$, le \textbf{signe de $\dH$} dit dans quel sens va la chaleur. On le lit sur un
\textbf{diagramme enthalpique} où l'on porte l'enthalpie $H$ des réactifs et des produits.

\begin{definitionb}[exo / endo / athermique]
$\dH = H_{\text{produits}} - H_{\text{réactifs}}$. Alors :
\begin{itemize}[leftmargin=1.5em,topsep=1pt,itemsep=1pt]
  \item \textbf{exothermique} : $\dH<0$, la réaction \emph{dégage} de la chaleur ; produits \emph{plus bas} que réactifs ;
  \item \textbf{endothermique} : $\dH>0$, la réaction \emph{absorbe} de la chaleur ; produits \emph{plus haut} ;
  \item \textbf{athermique} : $\dH\approx 0$, ni dégagement ni absorption.
\end{itemize}
\end{definitionb}

\begin{center}
\begin{tikzpicture}[scale=0.82]
% exo
\begin{scope}
  \draw[->,thick] (0,-0.2) -- (0,3.9) node[above,font=\small]{$H$};
  \draw[exocol,very thick] (0.4,3.1) -- (2.3,3.1); \node[anchor=south west,font=\scriptsize] at (0.45,3.15){Réactifs};
  \draw[exocol,very thick] (3.0,1.0) -- (4.9,1.0); \node[anchor=south west,font=\scriptsize] at (3.05,1.05){Produits};
  \draw[->,exocol,line width=1.1pt] (2.65,3.1) -- (2.65,1.0) node[midway,right,font=\scriptsize]{$\dH<0$};
  \node[exocol,font=\bfseries\scriptsize] at (2.6,3.65){exothermique};
\end{scope}
% endo
\begin{scope}[xshift=7.2cm]
  \draw[->,thick] (0,-0.2) -- (0,3.9) node[above,font=\small]{$H$};
  \draw[endocol,very thick] (0.4,1.0) -- (2.3,1.0); \node[anchor=south west,font=\scriptsize] at (0.45,1.05){Réactifs};
  \draw[endocol,very thick] (3.0,3.1) -- (4.9,3.1); \node[anchor=south west,font=\scriptsize] at (3.05,3.15){Produits};
  \draw[->,endocol,line width=1.1pt] (2.65,1.0) -- (2.65,3.1) node[midway,right,font=\scriptsize]{$\dH>0$};
  \node[endocol,font=\bfseries\scriptsize] at (2.6,3.65){endothermique};
\end{scope}
\end{tikzpicture}\\[2pt]
{\small\textbf{\textcolor{primary}{Figure}} : la flèche verticale donne immédiatement le signe de $\dH$.}
\end{center}

\begin{exempleb}[lire une équation]
Combustion de l'octane : \;$\ce{C8H18} + \tfrac{25}{2}\,\ce{O2} \rightarrow 8\,\ce{CO2} + 9\,\ce{H2O} + \text{Énergie}$.\;
Le « $+\,$Énergie » à droite signale une réaction \textbf{exothermique} ($\dH<0$) : les produits sont
\emph{moins} énergétiques que les réactifs. \emph{Toute combustion est exothermique.}
\end{exempleb}

\begin{attentionb}
Ne pas confondre \textbf{athermique} (une \emph{réaction} qui n'échange pas de chaleur, $\dH\approx0$) et
\textbf{adiabatique} (une \emph{enceinte} qui ne laisse passer aucune chaleur). Et « $\dH<0$ » signifie
que les produits ont \emph{moins} d'énergie : l'énergie est \emph{libérée} vers l'extérieur (pas « perdue »).
On ne manipule d'ailleurs jamais $H$ seul, seulement des \emph{variations} $\dH$ — point de départ du
Thème 2 (loi de Hess).
\end{attentionb}

\end{document}
